%!TEX program = lualatex
\documentclass[ngerman,headheight=70pt]{scrartcl}
\usepackage[ngerman]{babel}
\usepackage{blindtext}
\usepackage[utf8]{luainputenc}
\usepackage{scrlayer-scrpage}
\usepackage[T1]{fontenc}
\usepackage{fontspec}
\usepackage{multicol}
\usepackage{enumitem}
\usepackage{epigraph}
\usepackage[german=quotes,autostyle=true,strict=true]{csquotes}
\usepackage[bottom=2cm,footskip=8mm]{geometry}
\setmainfont{Source Sans Pro}
\MakeOuterQuote{"}

\setlength{\parindent}{0pt}
\setlength{\parskip}{5pt}

\pagestyle{plain.scrheadings}
\thispagestyle{scrheadings}
\lohead{\vspace*{-1.5cm}\includegraphics[scale=0.4]{../images/stupa_logo_rot.pdf}}
\rohead{\\Präsidium des\\Studierendenparlaments}

\lofoot[Universität Hamburg · Präsidium des Studierendenparlaments\\\hrule Von-Melle-Park 5 · D-20146 Hamburg]{Universität Hamburg · Präsidium des Studierendenparlaments\\\hrule Von-Melle-Park 5 · D-20146 Hamburg}
\cofoot[]{}
\cohead[\thepage]{}

\renewcommand{\thesubsection}{TOP \arabic{subsection}}
\renewcommand{\thesubsubsection}{\arabic{subsubsection}.}
\setcounter{subsection}{-1}

\begin{document}
    UHH · StuPa-Präsidium · Von-Melle-Park 5 · D-20146 Hamburg

    \section*{Beschluss des Studierendenparlamentes vom 26. Januar 2017}
    \subsection*{Wahlinformationen und -werbung auf Englisch verfügbar machen --
    Deutsch als Fremdsprache stärken}

    Das Präsidium des Studierendenparlaments wird dazu aufgefordert, die
    wichtigsten Wahlinformationen ab der Wahl zum Studierendenparlament 2018/19,
    welche im Dezember/Januar 2017/18 stattfinden wird, auch in Englisch zumindest
    auf der Website des Parlaments anzubieten. Den zum Studierendenparlament
    kandidierenden Listen wird beginnend mit der gleichen Wahl angeboten zumindest
    ihre kleinen Listendarstellungen auch in Englisch zumindest auf der Website
    des Parlaments veröffentlichen zu lassen.

    Ergänzend dazu wird die Universität Hamburg aufgefordert Angebote für Deutsch
    als Fremdsprache auszubauen, damit internationale Studierende eine bessere
    Möglichkeit bekommen, sich die Landessprache anzueignen, und in die Lage
    versetzt werden an der (Hochschul-)Politik teilzunehmen.

\end{document}
